%%% Ví dụ: Toán lớp 10 - Chủ đề Hàm số bậc hai %%%

%%% LOẠI 1: Trắc nghiệm %%%
%%%=============EX_1=============%%%
\begin{ex}%[0T3NB1-1]
	Parabol $y = x^2 - 4x + 3$ có đỉnh là điểm nào sau đây?
	\choice
	{$(2; -3)$}
	{$(-2; 1)$}
	{\True $(2; -1)$}
	{$(2; 1)$}
	\loigiai{
		Ta có $y = x^2 - 4x + 3 = (x-2)^2 - 1$.
		\\
		Đỉnh parabol là $I(2; -1)$.
	}
\end{ex}

%%% LOẠI 3: Trả lời ngắn %%%
%%%==============SA_1==============%%%
\begin{ex}%[0T3TH1-1]
	Cho hàm số $f(x) = 2x^2 - 8x + 5$. Tính giá trị nhỏ nhất của $f(x)$.
	\shortans{$-3$}
	\loigiai{
		Ta có $f(x) = 2(x^2 - 4x) + 5 = 2(x-2)^2 - 8 + 5 = 2(x-2)^2 - 3$.
		\\
		Vì $(x-2)^2 \geq 0$ nên $f(x) \geq -3$.
		\\
		Dấu bằng xảy ra khi $x = 2$. Vậy $\min f(x) = -3$.
	}
\end{ex}

%%% LOẠI 4: Tự luận %%%
%%%=============BT_1=============%%%
\begin{bt}%[0T3VD1-1]
	Cho parabol $(P): y = x^2 - 2x - 3$.
	\begin{enumerate}
		\item Tìm tọa độ đỉnh, trục đối xứng của $(P)$.
		\item Tìm giao điểm của $(P)$ với trục $Ox$.
		\item Vẽ đồ thị hàm số.
	\end{enumerate}
	\loigiai{
		\textbf{1.} Ta có $y = (x-1)^2 - 4$. Đỉnh $I(1; -4)$, trục đối xứng $x = 1$.

		\textbf{2.} Giao điểm với $Ox$: cho $y = 0$:
		\begin{eqnarray*}
			x^2 - 2x - 3 &=& 0 \\
			(x-3)(x+1) &=& 0 \\
			x = 3 &\text{hoặc}& x = -1
		\end{eqnarray*}
		Giao điểm: $A(-1; 0)$ và $B(3; 0)$.

		\textbf{3.} Bảng giá trị:
		\begin{center}
			\begin{tabular}{|c|c|c|c|c|c|}
				\hline
				$x$ & $-1$ & $0$ & $1$ & $2$ & $3$ \\
				\hline
				$y$ & $0$ & $-3$ & $-4$ & $-3$ & $0$ \\
				\hline
			\end{tabular}
		\end{center}
		Vẽ parabol qua các điểm trên với đỉnh $I(1; -4)$.
	}
\end{bt}
